% GENERATED FILE. Edit canonical lessons, then run npm run generate:books.
% canonical-source-hash: fnv1a64:01d919922a378b63
% canonical-lessons: ES-C577-argumentar, ES-C577-objecion, ES-C577-convencido, ES-C577-razonable, ES-C577-logico

\chapter{Defender una idea}
\label{ch:es-defender-una-idea}
\hlchaptermodality{577}

\begin{chapteropening}
\textbf{By the end of this chapter:} \emph{I can give reasons for an idea, answer an objection and say that it is reasonable.}

Complete the last lesson of chapter 577: I can give reasons for an idea, answer an objection and say that it is reasonable.
\end{chapteropening}

\section[argumentar]{argumentar — to argue, to give reasons}

\label{lesson:ES-C577-argumentar}



\begin{quote}
Before the new one: say the Spanish for the conviction, then the Spanish for the stance.
\end{quote}

\subsection*{You'll want to know: argumentar}
\textbf{argumentar} — \textquotedblleft{}to argue, to give reasons\textquotedblright{}.

\begin{quote}
Tienes que argumentar tu respuesta.
\end{quote}

\begin{quote}
\emph{You have to give reasons for your answer.}
\end{quote}

To argue a point with reasons, not to quarrel: a quarrel is \emph{discutir} or \emph{pelearse}.

\begin{cousinweb}
\emph{argumentar} is built on \emph{argumento}, a learned borrowing from Latin \emph{argumentum} 'proof', from \emph{arguere} 'to make clear, to prove' (Lewis \& Short).
\end{cousinweb}

\begin{grammarlens}[title={What you've built}]
One more word for telling, explaining and describing at B1.
\end{grammarlens}

\subsection*{Your turn}
\begin{itemize}
\raggedright
  \item \emph{Say it:} \emph{argumentar}
  \item \emph{Say it:} \emph{argumentar}, once more
  \item \emph{Say it:} a sentence of your own with \emph{argumentar}
  \item \emph{Recall:} say the Spanish for the conviction, then the Spanish for the stance, then say \emph{argumentar} again
\end{itemize}

\begin{tcolorbox}[breakable,colback=teal!4,colframe=teal!35!black,title={Before you move on}]
What does argumentar mean? (\textquotedblleft{}To argue, to give reasons\textquotedblright{}.) Say it once more.
\end{tcolorbox}
\section[la objeción]{la objeción — the objection}

\label{lesson:ES-C577-objecion}



\begin{quote}
Before the new one: say the Spanish for the stance, then the Spanish for to argue.
\end{quote}

\subsection*{You'll want to know: la objeción}
\textbf{la objeción} — \textquotedblleft{}the objection\textquotedblright{}.

\begin{quote}
¿Alguien tiene alguna objeción?
\end{quote}

\begin{quote}
\emph{Does anyone have any objection?}
\end{quote}

\emph{Poner una objeción} — to raise an objection.

\begin{cousinweb}
\emph{objeción} is a learned borrowing from Latin \emph{obiectio} 'a throwing against, an objection', built on \emph{obicere}, from \emph{ob} and \emph{iacere} 'to throw' (EtymDB).
\end{cousinweb}

\begin{grammarlens}[title={What you've built}]
One more word for telling, explaining and describing at B1.
\end{grammarlens}

\subsection*{Your turn}
\begin{itemize}
\raggedright
  \item \emph{Say it:} \emph{la objeción}
  \item \emph{Say it:} \emph{la objeción}, once more
  \item \emph{Say it:} a sentence of your own with \emph{la objeción}
  \item \emph{Recall:} say the Spanish for the stance, then the Spanish for to argue, then say \emph{la objeción} again
\end{itemize}

\begin{tcolorbox}[breakable,colback=teal!4,colframe=teal!35!black,title={Before you move on}]
What does la objeción mean? (\textquotedblleft{}The objection\textquotedblright{}.) Say it once more.
\end{tcolorbox}
\section[convencido]{convencido — convinced}

\label{lesson:ES-C577-convencido}



\begin{quote}
Before the new one: say the Spanish for to argue, then the Spanish for the objection.
\end{quote}

\subsection*{You'll want to know: convencido}
\textbf{convencido} — \textquotedblleft{}convinced\textquotedblright{}.

\begin{quote}
Estoy convencido de que es verdad.
\end{quote}

\begin{quote}
\emph{I am convinced that it is true.}
\end{quote}

With \emph{estar} and \emph{de que}: \emph{estoy convencida de que...}

\begin{cousinweb}
\emph{convencido} is built on \emph{convencer}, inherited from Latin \emph{convincere} (EtymDB) 'to overcome, to prove wrong', \emph{con} + \emph{vincere} 'to conquer' (Lewis \& Short).
\end{cousinweb}

\begin{grammarlens}[title={What you've built}]
One more word for telling, explaining and describing at B1.
\end{grammarlens}

\subsection*{Your turn}
\begin{itemize}
\raggedright
  \item \emph{Say it:} \emph{convencido}
  \item \emph{Say it:} \emph{convencido}, once more
  \item \emph{Say it:} a sentence of your own with \emph{convencido}
  \item \emph{Recall:} say the Spanish for to argue, then the Spanish for the objection, then say \emph{convencido} again
\end{itemize}

\begin{tcolorbox}[breakable,colback=teal!4,colframe=teal!35!black,title={Before you move on}]
What does convencido mean? (\textquotedblleft{}Convinced\textquotedblright{}.) Say it once more.
\end{tcolorbox}
\section[razonable]{razonable — reasonable}

\label{lesson:ES-C577-razonable}



\begin{quote}
Before the new one: say the Spanish for the objection, then the Spanish for convinced.
\end{quote}

\subsection*{You'll want to know: razonable}
\textbf{razonable} — \textquotedblleft{}reasonable\textquotedblright{}.

\begin{quote}
Es un precio razonable.
\end{quote}

\begin{quote}
\emph{It is a reasonable price.}
\end{quote}

The same for both genders: \emph{una idea razonable}.

\begin{cousinweb}
\emph{razonable} is inherited from Late Latin \emph{rationabilis}, built on \emph{ratio} 'reckoning, reason', the word that also gave \emph{razón} (EtymDB).
\end{cousinweb}

\begin{grammarlens}[title={What you've built}]
One more word for telling, explaining and describing at B1.
\end{grammarlens}

\subsection*{Your turn}
\begin{itemize}
\raggedright
  \item \emph{Say it:} \emph{razonable}
  \item \emph{Say it:} \emph{razonable}, once more
  \item \emph{Say it:} a sentence of your own with \emph{razonable}
  \item \emph{Recall:} say the Spanish for the objection, then the Spanish for convinced, then say \emph{razonable} again
\end{itemize}

\begin{tcolorbox}[breakable,colback=teal!4,colframe=teal!35!black,title={Before you move on}]
What does razonable mean? (\textquotedblleft{}Reasonable\textquotedblright{}.) Say it once more.
\end{tcolorbox}
\section[lógico]{lógico — logical; natural, to be expected}

\label{lesson:ES-C577-logico}



\begin{quote}
Before the new one: say the Spanish for convinced, then the Spanish for reasonable.
\end{quote}

\subsection*{You'll want to know: lógico}
\textbf{lógico} — \textquotedblleft{}logical; natural, to be expected\textquotedblright{}.

\begin{quote}
Es lógico que estés cansado.
\end{quote}

\begin{quote}
\emph{It is natural that you are tired.}
\end{quote}

\emph{Es lógico que} + subjunctive: it makes sense that...

\begin{cousinweb}
\emph{lógico} is a learned borrowing from Latin \emph{logicus}, taken from Greek \emph{logikos} 'of reasoning', built on \emph{logos} 'word, reason' (EtymDB).
\end{cousinweb}

\begin{grammarlens}[title={What you've built}]
One more word for telling, explaining and describing at B1.
\end{grammarlens}

\subsection*{Your turn}
\begin{itemize}
\raggedright
  \item \emph{Say it:} \emph{lógico}
  \item \emph{Say it:} \emph{lógico}, once more
  \item \emph{Say it:} a sentence of your own with \emph{lógico}
  \item \emph{Recall:} say the Spanish for convinced, then the Spanish for reasonable, then say \emph{lógico} again
\end{itemize}

\begin{tcolorbox}[breakable,colback=teal!4,colframe=teal!35!black,title={Before you move on}]
What does lógico mean? (\textquotedblleft{}Logical; natural, to be expected\textquotedblright{}.) Say it once more.
\end{tcolorbox}
